%% kapa: exam=1 round=29 year=2018 period=2 form=A subject=appraisal_law subject_ko=감정평가관계법규 seq=29 number=29
%% source: Q-Net 감정평가사 기출문제 (cst00309) artlSeq=5202808 file=extracted/1008/5202808/2018년 제29회 감정평가사 1차시험 기출문제/2018년 제29회 감정평가사 1차시험 기출문제/2교시 A형.pdf
%% conversion: kordoc 4.24.51 -> tools/gen1.py (auto); review=proofread-sonnet; flags=none
\begin{question}{29}
공간정보의 구축 및 관리 등에 관한 법령상 지목의 구분기준과 종류가 옳게 연결된 것은?
\begin{choices}
  \item 자연의 유수(流水)가 있거나 있을 것으로 예상되는 토지 - 하천
  \item 축산업 및 낙농업을 하기 위하여 초지를 조성한 토지 내의 주거용 건축물의 부지 - 목장용지
  \item 지하에서 용출되는 온수를 일정한 장소로 운송하는 송수관 및 저장시설의 부지 - 광천지
  \item 자동차 등의 판매 목적으로 설치된 물류장 - 주차장
  \item 아파트ㆍ공장 등 단일 용도의 일정한 단지 안에 설치된 통로 - 도로
\end{choices}
\end{question}
\answer{1}
